%!LW recipe=latexmk (lualatex)

\documentclass[a4paper,12pt]{article}

% =========================
% układ strony
% =========================
\usepackage[
  paper=a3paper,
  left=1.35cm,
  right=1.35cm,
  top=1.2cm,
  bottom=1.55cm,
  footskip=18pt
]{geometry}

\pagestyle{plain}
\raggedbottom

% =========================
% typografia i język
% =========================
\usepackage{microtype}
\usepackage{fontspec}
\usepackage{polyglossia}
\setdefaultlanguage{polish}
\setmainfont{STIX Two Text}

% =========================
% matematyka
% =========================
\usepackage{amsmath,amssymb,mathtools}
\usepackage{unicode-math}
\usepackage{diagbox}
\setmathfont{STIX Two Math}

% =========================
% dodatki
% =========================
\usepackage{xcolor}
\usepackage{multicol}
\usepackage{enumitem}
\usepackage[most]{tcolorbox}
\usepackage{booktabs} % tabelki
\usepackage{tikz}
\usepackage{csquotes}
\usetikzlibrary{automata, positioning, arrows.meta}
% =========================
% kolory
% =========================
\definecolor{TaskGray}{HTML}{F7F7F8}

% =========================
% globalna zwartość / odstępy
% =========================
\setlength{\parindent}{0pt}
\setlength{\parskip}{1.6pt}
\setlength{\columnsep}{0.9cm}
\setlength{\columnseprule}{0.5pt}
\linespread{1.02}

% =========================
% itemize, steps i podpunkty 
% =========================

\setlist[itemize]{
  label=\textbullet,
  leftmargin=1.2em,
  itemsep=0pt,
  topsep=0pt,
  parsep=0pt,
  partopsep=0pt
}

\newlist{steps}{enumerate}{1}
\setlist[steps]{
  label=\textbf{(\arabic*)},
  leftmargin=1.9em,
  itemsep=0.42em,
  topsep=0.2em,
  parsep=0pt,
  partopsep=0pt
}

\newlist{podpunkty}{enumerate}{1}
\setlist[podpunkty]{
  label=\textbf{(\alph*)},
  leftmargin=1.9em,
  itemsep=0.42em,
  topsep=0.2em,
  parsep=0pt,
  partopsep=0pt
}


% =========================
% indeks
% =========================
\newcommand{\indeks}[1]{%
  \par\bigskip
  \noindent
  {\setlength{\fboxsep}{8pt}%
   \setlength{\fboxrule}{1.3pt}%
   \fcolorbox{black}{white}{%
     \strut\fontsize{22}{26}\selectfont\bfseries #1%
   }}%
  \par\vspace{0.95em}
}

% =========================
% temat i bit
% =========================
\newcommand{\temat}[1]{%
  \par\medskip
  \noindent{\bfseries
      \fontsize{16}{20}\selectfont
      #1%
  }
  \par\medskip
}

\newcommand{\bit}[1]{%
  \par\medskip
  \noindent{\bfseries #1}%
  \par\smallskip
}

% =========================
% logbox
% =========================
\newtcolorbox{logbox}{
  enhanced,
  colback=TaskGray,
  colframe=white,
  boxrule=0.9pt,
  arc=0pt,
  left=6pt,
  right=6pt,
  top=5pt,
  bottom=5pt,
  boxsep=0pt,
  before skip=0.2em,
  after skip=0.7em
}

% =========================
% tw i def
% =========================
\newcommand{\twierdzenie}[2]{%
  \noindent
  #1
  \par\smallskip
  \noindent\hspace*{1.2em}$\Longrightarrow$
  \par\smallskip
  \noindent
  #2
  \par\smallskip
}

\newcommand{\definicja}[2]{\textbf{def #1}$\ \coloneqq\ $#2}

% =========================
% zadania (zwykłe i domowe)
% =========================
\newcommand{\zad}[1]{\textbf{z#1}}

\newcommand{\domek}[2]{%
  \tikz[baseline=(t.base)]{%
    \node[inner xsep=2.5pt, inner ysep=1.8pt, font=\bfseries] (t) {#2};%
    \draw[line width=0.6pt, line join=round, #1]
      (t.south west) -- (t.north west) --
      ([xshift=-1.5pt]t.north west) -- ([yshift=1.1ex]t.north) --
      ([xshift=1.5pt]t.north east) -- (t.north east) --
      (t.south east) -- cycle;%
  }%
}

\newcommand{\zd}[1]{\domek{}{z#1}}                              % domowe z numerem (pod kolokwium)
\newcommand{\zdc}{\domek{dash pattern=on 1.5pt off 1pt}{z}}    % domowe bez numeru (ciekawostka)
\newcommand{\itemzd}{\stepcounter{podpunktyi}\item[\domek{}{(\alph{podpunktyi})}]}   % domowy podpunkt w podpunkty
% =========================
% makra matematyczne
% =========================
\newcommand{\N}{\mathbb{N}}
\newcommand{\Z}{\mathbb{Z}}
\newcommand{\Q}{\mathbb{Q}}
\newcommand{\R}{\mathbb{R}}
\newcommand{\C}{\mathbb{C}}
\newcommand{\K}{\mathbb{K}}
\newcommand{\X}{\mathbb{X}}

\DeclareMathOperator{\Lin}{Lin}

\DeclarePairedDelimiter{\abs}{\lvert}{\rvert}
\DeclarePairedDelimiter{\norm}{\lVert}{\rVert}
\DeclarePairedDelimiter{\set}{\lbrace}{\rbrace}
\DeclarePairedDelimiter{\inner}{\langle}{\rangle}
 
\begin{document}

\bigskip
\bigskip
\indeks{ć1}

\bigskip
\textbf{def} \; \(X\) - zbiór, \; \(\mathcal{T} \subseteq 2^X\), \; \((X, \mathcal{T})\) - \textbf{przestrzeń topologiczna} \, \(\iff\)
\begin{enumerate}[label=\roman*), leftmargin=1.9em, itemsep=0.2em, topsep=0.2em, parsep=0pt, partopsep=0pt]
    \item \(\emptyset, X \in \mathcal{T}\)
    \item \(\forall_{\alpha \in A} \;\, U_\alpha \in \mathcal{T} \; \implies \displaystyle\bigcup_{\alpha \in A} U_\alpha \in \mathcal{T}\)
    \item \(A, B \in \mathcal{T} \implies A \cap B \in \mathcal{T}\)
\end{enumerate}

\smallskip
\textbf{def} \; elementy \(\mathcal{T}\) - \textbf{zbiory otwarte}, \; ich dopełnienia - \textbf{zbiory domknięte}

\bigskip
\begin{logbox}
\textbf{przykłady}
\begin{steps}
    \item \(X\) - zbiór, \; \(\mathcal{T} = 2^X\) \, - \textbf{topologia dyskretna}
    \item \(\mathcal{T} = \{\emptyset, X\}\) \, - \textbf{topologia antydyskretna}
\end{steps}
\end{logbox}

\bigskip
\textbf{pytanie} \; \(X = \mathbb{R}\) \, - czy przedziały otwarte są zbiorami otwartymi w topologii dyskretnej i antydyskretnej?

\smallskip
odp.: topologia dyskretna - tak, \; topologia antydyskretna - nie


\bigskip
\zad{1} \; czy \((X, \mathcal{T})\) - przestrzeń topologiczna?
\begin{podpunkty}
    \item \(X = [0, \infty), \; \mathcal{T} = \{\emptyset, [0, \infty), (a, \infty) \text{ dla } a \geq 0\}\)
    \begin{enumerate}[label=\roman*), leftmargin=1.9em, itemsep=0.2em, topsep=0.2em, parsep=0pt, partopsep=0pt]
        \item trywialne \, \(\checkmark\)
        \item \(U_\alpha \coloneqq (a_\alpha, \infty), \; a_\alpha \geq 0, \; \alpha \in A \implies \bigcup_{\alpha \in A} U_\alpha = \bigcup_{\alpha \in A} (a_\alpha, \infty) \overset{(*)}{=} \big(\inf_{\alpha \in A} a_\alpha, \infty\big) \in \mathcal{T}\) \, \(\checkmark\)

        \smallskip
        \((*)\):
        \begin{itemize}
            \item \(\subseteq\): \, \(x \in \bigcup_{\alpha \in A} (a_\alpha, \infty) \iff \exists_{\alpha_0 \in A} \;\, x \in (a_{\alpha_0}, \infty) \implies x \in \big(\inf_{\alpha \in A} a_\alpha, \infty\big)\), \, bo \(\forall_{\alpha \in A} \;\, a_\alpha \geq \inf_{\alpha \in A} a_\alpha \implies (a_\alpha, \infty) \subseteq \big(\inf_{\alpha \in A} a_\alpha, \infty\big)\)
            \item \(\supseteq\): \, \(x \in \big(\inf_{\alpha \in A} a_\alpha, \infty\big) \overset{\text{z def kresu dolnego}}{\implies} \exists_{\alpha_0 \in A} \;\, a_{\alpha_0} < x \implies x \in (a_{\alpha_0}, \infty) \implies x \in \bigcup_{\alpha \in A} (a_\alpha, \infty)\)
        \end{itemize}
        \item \(\alpha_1, \alpha_2 \in A \implies (a_{\alpha_1}, \infty) \cap (a_{\alpha_2}, \infty) = \big(\max\{a_{\alpha_1}, a_{\alpha_2}\}, \infty\big) \in \mathcal{T}\) \, \(\checkmark\)
    \end{enumerate}
    \item \(X = [0, \infty), \; \mathcal{T} = \{\emptyset, [0, \infty), [a, \infty) \text{ dla } a \geq 0\}\)
    \begin{enumerate}[label=\roman*), leftmargin=1.9em, itemsep=0.2em, topsep=0.2em, parsep=0pt, partopsep=0pt]
        \setcounter{enumi}{1}
        \item \(\bigcup_{n=1}^{\infty} \big[\frac{1}{n}, \infty\big) = (0, \infty) \notin \mathcal{T}\) \, \(\times\)
    \end{enumerate}
    \item \(X = \mathbb{R}, \; U \in \mathcal{T} \iff U\) - suma przedziałów otwartych \(\iff \exists_{\Gamma \text{ - zbiór}} \; \exists_{a_\gamma, b_\gamma \in \overline{\mathbb{R}}, \; \gamma \in \Gamma} \;\, U = \bigcup_{\gamma \in \Gamma} (a_\gamma, b_\gamma)\)
    \begin{enumerate}[label=\roman*), leftmargin=1.9em, itemsep=0.2em, topsep=0.2em, parsep=0pt, partopsep=0pt]
        \item trywialne \, \(\checkmark\)
        \item \(U_\alpha = \bigcup_{\gamma_\alpha \in \Gamma_\alpha} (a_{\gamma_\alpha}, b_{\gamma_\alpha}) \implies \bigcup_{\alpha \in A} U_\alpha = \bigcup_{\alpha \in A} \bigcup_{\gamma_\alpha \in \Gamma_\alpha} (a_{\gamma_\alpha}, b_{\gamma_\alpha}) \in \mathcal{T}\) \, (suma sum przedziałów otwartych jest sumą przedziałów otwartych) \, \(\checkmark\)
        \item \(A = \bigcup_{\gamma_1 \in \Gamma_1} (a_{\gamma_1}, b_{\gamma_1}), \; B = \bigcup_{\gamma_2 \in \Gamma_2} (a_{\gamma_2}, b_{\gamma_2}) \implies\)
        \begin{align*}
        A \cap B &= \bigcup_{\gamma_1 \in \Gamma_1} (a_{\gamma_1}, b_{\gamma_1}) \cap \bigcup_{\gamma_2 \in \Gamma_2} (a_{\gamma_2}, b_{\gamma_2}) \overset{(*)}{=} \bigcup_{\gamma_1 \in \Gamma_1} \bigcup_{\gamma_2 \in \Gamma_2} \big((a_{\gamma_1}, b_{\gamma_1}) \cap (a_{\gamma_2}, b_{\gamma_2})\big) = \\
        &= \bigcup_{\gamma_1 \in \Gamma_1} \bigcup_{\gamma_2 \in \Gamma_2} \underbrace{\big(\max\{a_{\gamma_1}, a_{\gamma_2}\}, \min\{b_{\gamma_1}, b_{\gamma_2}\}\big)}_{= \, \emptyset \text{ dla } \max \geq \min} \in \mathcal{T} \quad \checkmark
        \end{align*}
        \((*)\) rozdzielność przecięcia względem sumy \(A \cap (B \cup C) = (A \cap B) \cup (A \cap C)\), zastosowana dwukrotnie (raz dla każdej sumy)
    \end{enumerate}
\end{podpunkty}

\bigskip
\textbf{def} \; topologia z \zad{1}(c) \, - \textbf{topologia standardowa} na \(\mathbb{R}\)

\smallskip
\textbf{wniosek} \; w top. st. na \(\mathbb{R}\) przedziały otwarte są zbiorami otwartymi, a przedziały domknięte domkniętymi (\([a, b] = \big((-\infty, a) \cup (b, \infty)\big)^c\))

\bigskip
\textbf{def} \; \(d: X \times X \to [0, \infty)\) - \textbf{metryka} \, \(\iff \, \forall_{x, y, z \in X}\)
\begin{enumerate}[label=\roman*), leftmargin=1.9em, itemsep=0.2em, topsep=0.2em, parsep=0pt, partopsep=0pt]
    \item \(d(x, y) = 0 \iff x = y\)
    \item \(d(x, y) = d(y, x)\)
    \item \(d(x, y) \leq d(x, z) + d(z, y)\)
\end{enumerate}

\bigskip
niech \((X, d)\) - przestrzeń metryczna

\smallskip
\textbf{fakt} \; metryka zadaje topologię \(\mathcal{T}_d\): \, \(U \in \mathcal{T}_d \iff \forall_{x \in U} \; \exists_{\varepsilon > 0} \;\, B(x, \varepsilon) \subseteq U\)

\bigskip
\textbf{def} \; \((X, \mathcal{T})\) - przestrzeń topologiczna, \; \(\mathcal{B} \subseteq \mathcal{T}\) - \textbf{baza topologii} \, \(\iff\) każdy \(U \in \mathcal{T}\) jest sumą pewnej podrodziny z \(\mathcal{B}\)

\bigskip
\begin{logbox}
\textbf{przykład} \; w topologii standardowej na \(\mathbb{R}\) bazą są przedziały otwarte, a także przedziały otwarte o końcach wymiernych \(\{(a, b) : a, b \in \mathbb{Q}\}\)
\end{logbox}

\smallskip
\textbf{wniosek} \; \(\mathbb{R}\) ma przeliczalną bazę

\bigskip
\textbf{pytanie} \; czy standardowa metryka na \(\mathbb{R}\) zadaje topologię standardową na \(\mathbb{R}\)?

\smallskip
\(d(x, y) = |x - y|\), \; \(\mathcal{T}_d\) - topologia zadana przez metrykę, \; \(\mathcal{T}\) - topologia standardowa na \(\mathbb{R}\), \; czy \(\mathcal{T}_d = \mathcal{T}\)?
\begin{itemize}
    \item \(\subseteq\): \, \(U \in \mathcal{T}_d \overset{\text{def}}{\iff} \forall_{x \in U} \; \exists_{\varepsilon_x > 0} \;\, B(x, \varepsilon_x) \subseteq U \implies U \subseteq \bigcup_{x \in U} B(x, \varepsilon_x) \subseteq U\), \, czyli \(U = \bigcup_{x \in U} B(x, \varepsilon_x) \in \mathcal{T}\), \, bo \(B(x, \varepsilon) = (x - \varepsilon, x + \varepsilon)\)
    \item \(\supseteq\): \, \(U \in \mathcal{T} \implies U = \bigcup_{\alpha \in A} B(x_\alpha, \varepsilon_\alpha) \implies \forall_{x \in U} \; \exists_{\alpha_0 \in A} \;\, x \in B(x_{\alpha_0}, \varepsilon_{\alpha_0}) \subseteq \bigcup_{\alpha \in A} B(x_\alpha, \varepsilon_\alpha) = U\) \\
    \(\varepsilon_x \coloneqq \varepsilon_{\alpha_0} - d(x, x_{\alpha_0}) > 0 \overset{\text{z nierówności trójkąta}}{\implies} B(x, \varepsilon_x) \subseteq B(x_{\alpha_0}, \varepsilon_{\alpha_0}) \subseteq U \implies U \in \mathcal{T}_d\)
\end{itemize}

\smallskip
odp.: tak

\newpage
\zad{2} \; na \(\mathbb{R}^2\) definiujemy:
\smallskip
\begin{itemize}
    \item \(d_1(x, y) = |x_1 - y_1| + |x_2 - y_2|\)
    \item \(d_2(x, y) = \big(|x_1 - y_1|^2 + |x_2 - y_2|^2\big)^{\frac{1}{2}}\)
    \smallskip
    \item \(d_\infty(x, y) = \max\{|x_1 - y_1|, |x_2 - y_2|\}\)
\end{itemize}
\smallskip
sprawdź, czy to metryki na \(\mathbb{R}^2\)
\smallskip
\begin{podpunkty}
    \item \(d_1\):
    \begin{enumerate}[label=\roman*), leftmargin=1.9em, itemsep=0.2em, topsep=0.2em, parsep=0pt, partopsep=0pt]
        \item \(d_1(x, y) = 0 \overset{|\cdot| \geq 0}{\iff} x_1 = y_1 \land x_2 = y_2 \iff x = y\) \, \(\checkmark\)
        \item trywialne z \(|a - b| = |b - a|\) \, \(\checkmark\)
        \item
        \begin{align*}
        d_1(x, y) &= |x_1 - y_1| + |x_2 - y_2| = |(x_1 - z_1) + (z_1 - y_1)| + |(x_2 - z_2) + (z_2 - y_2)| \overset{\text{z nierówności trójkąta}}{\leq} \\
        &\leq |x_1 - z_1| + |x_2 - z_2| + |z_1 - y_1| + |z_2 - y_2| = d_1(x, z) + d_1(z, y) \quad \checkmark
        \end{align*}
    \end{enumerate}
    \item \(d_2\):
    \begin{enumerate}[label=\roman*), leftmargin=1.9em, itemsep=0.2em, topsep=0.2em, parsep=0pt, partopsep=0pt]
        \item trywialne \, \(\checkmark\)
        \item trywialne \, \(\checkmark\)
        \item ozn.: \(a = x - z, \; b = z - y\) \, \((a_i = x_i - z_i, \; b_i = z_i - y_i)\), \; wykażemy \(d_2(x, y) \leq d_2(x, z) + d_2(z, y)\)
        \begin{align*}
        & \big(|x_1 - y_1|^2 + |x_2 - y_2|^2\big)^{\frac{1}{2}} \leq d_2(x, z) + d_2(z, y) \iff \\
        \iff \; & \big(|a_1 + b_1|^2 + |a_2 + b_2|^2\big)^{\frac{1}{2}} \leq \big(a_1^2 + a_2^2\big)^{\frac{1}{2}} + \big(b_1^2 + b_2^2\big)^{\frac{1}{2}} \qquad \big/ \, (\cdot)^2 \; \text{(obie strony nieujemne)} \\
        \iff \; & (a_1 + b_1)^2 + (a_2 + b_2)^2 \leq a_1^2 + a_2^2 + b_1^2 + b_2^2 + 2\sqrt{a_1^2 + a_2^2} \, \sqrt{b_1^2 + b_2^2} \\
        \iff \; & a_1^2 + b_1^2 + a_2^2 + b_2^2 + 2a_1b_1 + 2a_2b_2 \leq a_1^2 + a_2^2 + b_1^2 + b_2^2 + 2\sqrt{a_1^2 + a_2^2} \, \sqrt{b_1^2 + b_2^2} \\
        \overset{\text{po redukcji}}{\iff} \; & a_1b_1 + a_2b_2 \leq \sqrt{a_1^2 + a_2^2} \, \sqrt{b_1^2 + b_2^2}
        \end{align*}
        prawdziwe, bo \(a_1b_1 + a_2b_2 \leq |a_1b_1 + a_2b_2| \overset{\text{z nierówności Cauchy'ego-Schwarza}}{\leq} \sqrt{a_1^2 + a_2^2} \, \sqrt{b_1^2 + b_2^2}\) \, \(\checkmark\)
    \end{enumerate}
    \item \(d_\infty\):
    \begin{enumerate}[label=\roman*), leftmargin=1.9em, itemsep=0.2em, topsep=0.2em, parsep=0pt, partopsep=0pt]
        \item trywialne \, \(\checkmark\)
        \item trywialne \, \(\checkmark\)
        \item \((*)\) \; \(\max\{u, v\} \leq M \iff u \leq M \land v \leq M\)
        \begin{align*}
        |x_1 - y_1| &= |x_1 - z_1 + z_1 - y_1| \overset{\text{z nierówności trójkąta}}{\leq} |x_1 - z_1| + |z_1 - y_1| \leq \\
        &\leq \max\{|x_1 - z_1|, |x_2 - z_2|\} + \max\{|z_1 - y_1|, |z_2 - y_2|\} \eqqcolon M
        \end{align*}
        analogicznie dla \(i = 2\): \, \(|x_2 - y_2| \leq M\), \; zatem
        \begin{align*}
        d_\infty(x, y) = \max\{|x_1 - y_1|, |x_2 - y_2|\} \overset{(*)}{\leq} \max\{|x_1 - z_1|, |x_2 - z_2|\} + \max\{|z_1 - y_1|, |z_2 - y_2|\} = d_\infty(x, z) + d_\infty(z, y) \quad \checkmark
        \end{align*}
    \end{enumerate}
\end{podpunkty}


\bigskip
\bigskip
\indeks{ć2}

\bigskip
\textbf{pytanie} \; jakie są inkluzje między topologiami \(\mathcal{T}_{d_1}, \mathcal{T}_{d_2}, \mathcal{T}_{d_\infty}\)?

\smallskip
będziemy rozważać kule:
\begin{itemize}
    \item \(B_{d_2}(0, 1) = \{x \in \mathbb{R}^2 : d_2(0, x) < 1\} = \{x \in \mathbb{R}^2 : \big(|x_1|^2 + |x_2|^2\big)^{\frac{1}{2}} < 1\}\) \, - koło
    \item \(B_{d_1}(0, 1) = \{x \in \mathbb{R}^2 : d_1(0, x) < 1\} = \{x \in \mathbb{R}^2 : |x_1| + |x_2| < 1\}\) \, - romb
    \item \(B_{d_\infty}(0, 1) = \{x \in \mathbb{R}^2 : d_\infty(0, x) < 1\} = \{x \in \mathbb{R}^2 : \max\{|x_1|, |x_2|\} < 1\}\) \, - kwadrat
\end{itemize}

\smallskip
\(\mathcal{T}_{d_1} = \mathcal{T}_{d_2} = \mathcal{T}_{d_\infty}\), \, bo
\begin{itemize}
    \item \(U \in \mathcal{T}_{d_2} \iff \forall_{x \in U} \; \exists_{\varepsilon > 0} \;\, B_{d_2}(x, \varepsilon) \subseteq U\)
    \item \(U \in \mathcal{T}_{d_1} \iff \forall_{x \in U} \; \exists_{\varepsilon > 0} \;\, B_{d_1}(x, \varepsilon) \subseteq U\)
\end{itemize}

\smallskip
a w każdej kuli jednej metryki o środku \(x\) mieści się mniejsza kula każdej innej o tym samym środku:
\begin{align*}
d_\infty \leq d_2 \leq d_1 \leq 2 \, d_\infty \implies B_{d_\infty}\big(x, \tfrac{\varepsilon}{2}\big) \subseteq B_{d_1}(x, \varepsilon) \subseteq B_{d_2}(x, \varepsilon) \subseteq B_{d_\infty}(x, \varepsilon)
\end{align*}

\bigskip
\textbf{def} \; \((X, \mathcal{T})\) - przestrzeń topologiczna, \; \(A \subseteq X\), \; \textbf{wnętrze} \(A\): \, \(\operatorname{int} A \coloneqq \bigcup \{U \subseteq A : U \in \mathcal{T}\}\) \, (największy zbiór otwarty zawarty w \(A\))

\bigskip
\zad{3} \; na \((\mathbb{R}, d_e)\) wyznacz
\begin{podpunkty}
    \item \(\operatorname{int}[a, b] \overset{?}{=} (a, b)\)
    \begin{align*}
    x_0 \in \operatorname{int}[a, b] \overset{\text{def wnętrza}}{\iff} \exists_{U_0 \subseteq [a, b], \; U_0 \in \mathcal{T}_{\text{st}}} \;\, x_0 \in U_0 \overset{(1)}{\iff} \exists_{r > 0} \;\, B(x_0, r) \subseteq [a, b] \overset{(2)}{\iff} x_0 \in (a, b)
    \end{align*}
    \begin{steps}
        \item \(\implies\): z def zbioru otwartego w przestrzeni metrycznej, \; \(\impliedby\): kule otwarte są zbiorami otwartymi, bierzemy \(U_0 = B(x_0, r)\)
        \item \(\implies\): \((x_0 - r, x_0 + r) \subseteq [a, b] \implies a \leq x_0 - r < x_0 < x_0 + r \leq b\), \; \(\impliedby\): \(r = \min\{x_0 - a, b - x_0\} > 0\), wtedy \((x_0 - r, x_0 + r) \subseteq (a, b) \subseteq [a, b]\)
    \end{steps}
    \item \(\operatorname{int}[a, b) \overset{?}{=} (a, b)\)

    \smallskip
    \textbf{monotoniczność wnętrza} \; \(A \subseteq B \implies \operatorname{int} A \subseteq \operatorname{int} B\), \, bo \(U \subseteq A \subseteq B \implies \bigcup \{U \subseteq A : U \in \mathcal{T}\} \subseteq \bigcup \{U \subseteq B : U \in \mathcal{T}\}\)
    \begin{itemize}
        \item \(\subseteq\): \, \([a, b) \subseteq [a, b] \overset{\text{z monotoniczności}}{\implies} \operatorname{int}[a, b) \subseteq \operatorname{int}[a, b] \overset{\text{(a)}}{=} (a, b)\)
        \item \(\supseteq\): \, \((a, b) \in \mathcal{T}_{\text{st}}, \; (a, b) \subseteq [a, b) \implies (a, b) \subseteq \bigcup \{U \subseteq [a, b) : U \in \mathcal{T}_{\text{st}}\} = \operatorname{int}[a, b)\)
    \end{itemize}
    \item \(\operatorname{int}(a, b) \overset{?}{=} (a, b)\)
    \begin{itemize}
        \item \(\subseteq\): \, \(\forall_A \;\, \operatorname{int} A \subseteq A\)
        \item \(\supseteq\): \, \((a, b) \in \mathcal{T}_{\text{st}}, \; (a, b) \subseteq (a, b) \implies (a, b) \subseteq \bigcup \{U \subseteq (a, b) : U \in \mathcal{T}_{\text{st}}\} = \operatorname{int}(a, b)\)
    \end{itemize}
\end{podpunkty}

\smallskip
\textbf{wniosek} \; \(U \in \mathcal{T} \iff \operatorname{int} U = U\)

\bigskip
\textbf{def} \; \textbf{domknięcie} \, \(\overline{A} \coloneqq \bigcap \{B \supseteq A : B^c \in \mathcal{T}\}\) \, (najmniejszy zbiór domknięty zawierający \(A\))

\smallskip
\textbf{fakt} \; \(A = \overline{A} \iff A\) - domknięty

\bigskip
\begin{align*}
\overline{A}^{\,c} \overset{\text{z praw de Morgana}}{=} \bigcup \{B^c : B \supseteq A, \; B^c \in \mathcal{T}\} \overset{B \supseteq A \, \iff \, B^c \subseteq A^c}{=} \bigcup \{B^c \subseteq A^c : B^c \in \mathcal{T}\} = \bigcup \{U \subseteq A^c : U \in \mathcal{T}\} \overset{\text{def int}}{=} \operatorname{int} A^c
\end{align*}

\smallskip
\textbf{wniosek} \; \(\overline{A} = \big(\operatorname{int} A^c\big)^c\)

\bigskip
\zad{4} \; w \((\mathbb{R}, d)\), \; \(d(x, y) = |x - y|\) \, wyznacz
\begin{podpunkty}
    \item \(\overline{(a, b)}\)

    \smallskip
    \((a, b)^c = (-\infty, a] \cup [b, \infty)\)
    \begin{align*}
    \overline{(a, b)} \overset{\overline{A} = (\operatorname{int} A^c)^c}{=} \big(\operatorname{int}\big((-\infty, a] \cup [b, \infty)\big)\big)^c = \big((-\infty, a) \cup (b, \infty)\big)^c = [a, b]
    \end{align*}
    po co tak? \, wnętrze można badać punktowo, domknięcie z definicji nie
    \item \(\overline{[a, b]} \overset{?}{=} [a, b]\)

    \smallskip
    \([a, b]^c = (-\infty, a) \cup (b, \infty) \in \mathcal{T}_{\text{st}} \implies [a, b]\) - domknięty \(\overset{\text{fakt}}{\implies} \overline{[a, b]} = [a, b]\)
    \item \(\overline{[a, b)} \overset{?}{=} [a, b]\)
    \begin{align*}
    [a, b] \overset{\text{(a)}}{=} \overline{(a, b)} \overset{\text{z monotoniczności domknięcia}}{\subseteq} \overline{[a, b)} \subseteq \overline{[a, b]} \overset{\text{(b)}}{=} [a, b]
    \end{align*}
\end{podpunkty}

\bigskip
\textbf{def} \; \((X, \mathcal{T})\) - przestrzeń topologiczna, \; \(A \subseteq X\), \; \textbf{brzeg} \(A \eqqcolon \partial A\): \, \(x \in \partial A \iff \forall_{U \in \mathcal{T}, \; x \in U} \;\, U \cap A \neq \emptyset \land U \cap A^c \neq \emptyset\)

\smallskip
\textbf{fakt} \; \(\partial A = \overline{A} \setminus \operatorname{int} A\)

\smallskip
\textbf{wniosek} \; \(\partial A\) - domknięty

\bigskip
\textbf{tw o trójpodziale} \; \((X, \mathcal{T})\) - przestrzeń topologiczna, \; \(A \subseteq X \implies X = \operatorname{int} A \sqcup \operatorname{int} A^c \sqcup \partial A\)

\bigskip
\begin{logbox}
\textbf{przykład} \; \(A = \mathbb{Q} \subseteq (\mathbb{R}, \mathcal{T}_{\text{st}})\) \, - \(\operatorname{int} \mathbb{Q} = \emptyset, \;\; \operatorname{int} \mathbb{Q}^c = \emptyset, \;\; \partial \mathbb{Q} = \mathbb{R}\)
\end{logbox}

\bigskip
\textbf{def} \; \((X, \mathcal{T})\) - przestrzeń topologiczna, \; \(Y \subseteq X\), \; \((Y, \mathcal{T}_Y)\) - przestrzeń topologiczna z \textbf{topologią indukowaną} z \(X\), \, gdzie \(\mathcal{T}_Y \coloneqq \{Y \cap U : U \in \mathcal{T}\}\)

\bigskip
\begin{logbox}
\textbf{przykład} \; \(X = \mathbb{R}, \; Y = [0, \infty)\), \, wtedy \([0, 1) \in \mathcal{T}_Y\), \, bo \([0, 1) = \underbrace{[0, \infty)}_{= \, Y} \cap \underbrace{(-1, 1)}_{\in \, \mathcal{T}_X}\)
\end{logbox}

\bigskip
\textbf{def} \; \((X, \mathcal{T})\) - przestrzeń topologiczna, \; \((x_n)_{n=1}^{\infty} \subset X\), \; \textbf{zbiór graniczny} \, \(\operatorname{Lim}(x_n) \coloneqq \{x_0 \in X : \forall_{U \in \mathcal{T}, \; x_0 \in U} \;\, \#\{n \in \mathbb{N} : x_n \notin U\} < \infty\}\)

\bigskip
\begin{logbox}
\textbf{przykład}
\begin{podpunkty}
    \item \((X, \mathcal{T})\) - przestrzeń topologiczna z topologią antydyskretną, \; \((x_n) \subset X\), \; \(\operatorname{Lim}(x_n) = \, ?\)

    \smallskip
    \(\forall_{x_0} \;\, x_0 \notin \emptyset\) \, (sprawdzamy tylko \(U = X\))

    \(\forall_{x_0} \;\, x_0 \in X\) \, i \, \(\#\{n \in \mathbb{N} : x_n \notin X\} = 0 < \infty\)

    więc \(\operatorname{Lim}(x_n) = X\)
    \item \(\mathcal{T}\) - topologia dyskretna, \; \(\operatorname{Lim}(x_n) = \, ?\)

    \smallskip
    intuicyjnie: do sprawdzenia jest każdy \(U\), w tym singleton \(\{x_0\}\), więc zbiór graniczny będzie mały lub pusty

    \smallskip
    ustalmy \(x_0 \in X\): \, \(x_0 \in \operatorname{Lim}(x_n) \iff (x_n)\) jest od pewnego momentu stale równy \(x_0\) \, \(\big(\#\{n \in \mathbb{N} : x_n \neq x_0\} < \infty\big)\)
    \begin{itemize}
        \item \(\implies\): \, \(x_0 \in \operatorname{Lim}(x_n) \implies \forall_{U \in \mathcal{T}, \; x_0 \in U} \;\, \#\{n \in \mathbb{N} : x_n \notin U\} < \infty \;\, \land \;\, \{x_0\} \in \mathcal{T} \implies \#\{n : x_n \notin \{x_0\}\} < \infty\)
        \item \(\impliedby\): \, \(\#\{n : x_n \neq x_0\} < \infty \;\, \land \;\, \forall_{U \in \mathcal{T}, \; x_0 \in U} \;\, \{x_0\} \subseteq U \implies \forall_{U \in \mathcal{T}, \; x_0 \in U} \;\, \#\{n : x_n \notin U\} \leq \#\{n : x_n \neq x_0\} < \infty \overset{\text{def}}{\iff} x_0 \in \operatorname{Lim}(x_n)\)
    \end{itemize}
    \smallskip
    ciąg może być od pewnego momentu stały tylko dla jednej wartości (trywialne), więc
    \[
    \operatorname{Lim}(x_n) = \begin{cases} \{x_0\}, & \exists_{x_0} \;\, \#\{n \in \mathbb{N} : x_n \neq x_0\} < \infty \\ \emptyset, & \text{w przeciwnym przypadku} \end{cases}
    \]
\end{podpunkty}
\end{logbox}

\bigskip
\textbf{def} \; \((X, \mathcal{T})\) - przestrzeń topologiczna, \; \((X, \mathcal{T})\) - \textbf{przestrzeń Hausdorffa} \, \(\iff \, \forall_{x_1, x_2 \in X, \; x_1 \neq x_2} \; \exists_{U_1, U_2 \in \mathcal{T}} \;\, x_1 \in U_1 \land x_2 \in U_2 \land U_1 \cap U_2 = \emptyset\)

\bigskip
\begin{logbox}
\textbf{przykład} \; \(\mathbb{R}^n, \; \mathbb{C}^n\) \, - przestrzenie Hausdorffa
\end{logbox}

\newpage
\zdc \; pokazać, że w przestrzeni Hausdorffa \(\#\operatorname{Lim}(x_n) \leq 1\)
\smallskip
\begin{itemize}
    \item przypuśćmy, że \(\exists_{x_0, y_0} \;\, x_0 \neq y_0 \land x_0, y_0 \in \operatorname{Lim}(x_n)\)
    \item z def przestrzeni Hausdorffa \(\exists_{U_{x_0}, U_{y_0} \in \mathcal{T}} \;\, x_0 \in U_{x_0} \land y_0 \in U_{y_0} \land U_{x_0} \cap U_{y_0} = \emptyset\)
    \item \(x_0 \in \operatorname{Lim}(x_n)\), zatem \(\#\{n \in \mathbb{N} : x_n \notin U_{x_0}\} < \infty\), \; \(\{n : x_n \notin U_{x_0}\} \supseteq \{n : x_n \in U_{y_0}\}\) \, (bo \(U_{x_0} \cap U_{y_0} = \emptyset\))
    \item zatem \(\#\{n : x_n \in U_{y_0}\} \leq \#\{n : x_n \notin U_{x_0}\} < \infty\)
    \item zatem \(\#\{n : x_n \notin U_{y_0}\} = \infty\), \, bo \(\mathbb{N} = \{n : x_n \in U_{y_0}\} \sqcup \{n : x_n \notin U_{y_0}\}, \; \#\mathbb{N} = \infty\) \, - sprzeczność z założeniem, że \(y_0 \in \operatorname{Lim}(x_n)\)
    \item więc \(\nexists_{x_0 \neq y_0} \;\, x_0, y_0 \in \operatorname{Lim}(x_n)\), więc \(\#\operatorname{Lim}(x_n) \leq 1\) \, \(\blacksquare\)
\end{itemize}

\bigskip
\zdc \; \(X\) - przestrzeń spójna \(\implies\) jedynymi zbiorami otwarto-domkniętymi są \(X\) i \(\emptyset\)
\begin{itemize}
    \item \(X\) i \(\emptyset\) są zawsze otwarto-domknięte z i) \(\land \; X^c = \emptyset, \; \emptyset^c = X\)
    \item \(X\) - spójna \(\iff \nexists_{A, B \text{ - otwarte, rozłączne, niepuste}} \;\, X = A \sqcup B\)
    \item przypuśćmy, że \(\exists_{A \subseteq X, \; A \notin \{\emptyset, X\}}\) \, \(A\) - otwarto-domknięty
    \item wtedy \(A \in \mathcal{T}\) z założenia
    \item \(A^c \in \mathcal{T}\), bo \(A\) - domknięty
    \item \(A, A^c\) - otwarte, rozłączne i niepuste \(\land \; X = A \sqcup A^c\) \, - sprzeczność ze spójnością \, \(\blacksquare\)
\end{itemize}


\bigskip
\bigskip
\indeks{ć3}

\bigskip
\zd{5} \; \((X, d)\) - przestrzeń metryczna, \; \(\tilde{d}(x, y) = \frac{d(x, y)}{1 + d(x, y)}\)

\smallskip
\((*)\) \; \(\tilde{d} \in [0, 1)\) \, - to pozwala normować każdą metrykę \(d\)
\begin{podpunkty}
    \item pokazać, że \(\tilde{d}\) - metryka na \(X\)
    \begin{enumerate}[label=\roman*), leftmargin=1.9em, itemsep=0.2em, topsep=0.2em, parsep=0pt, partopsep=0pt]
        \item trywialne
        \item trywialne
        \item
        \begin{align*}
        & \tilde{d}(x, y) \leq \tilde{d}(x, z) + \tilde{d}(z, y) \iff \\
        \iff \; & \frac{d(x, y)}{1 + d(x, y)} \leq \frac{d(x, z)}{1 + d(x, z)} + \frac{d(z, y)}{1 + d(z, y)} \\
        \overset{(1)}{\impliedby} \; & \frac{d(x, y)}{1 + d(x, y)} \leq \frac{d(x, z) + d(z, y)}{1 + d(x, z) + d(z, y)} \\
        \overset{(2)}{\iff} \; & d(x, y) \leq d(x, z) + d(z, y) \qquad \text{(nierówność trójkąta dla } d\text{)} \quad \blacksquare
        \end{align*}
        \begin{steps}
            \item \(\frac{d(x, z) + d(z, y)}{1 + d(x, z) + d(z, y)} = \frac{d(x, z)}{1 + d(x, z) + d(z, y)} + \frac{d(z, y)}{1 + d(x, z) + d(z, y)} \overset{d \geq 0}{\leq} \frac{d(x, z)}{1 + d(x, z)} + \frac{d(z, y)}{1 + d(z, y)}\)
            \item \(\forall_{a, b \geq 0} \;\, \frac{a}{1 + a} \leq \frac{b}{1 + b} \iff a(1 + b) \leq b(1 + a) \iff a + ab \leq b + ab \iff a \leq b\)
        \end{steps}
    \end{enumerate}
    \item czy \(\mathcal{T}_d = \mathcal{T}_{\tilde{d}}\)? \, tak

    \smallskip
    \(\varphi(t) = \frac{t}{1 + t}\) - ściśle rosnąca bijekcja \([0, \infty) \to [0, 1)\), \; \(\varphi^{-1}(s) = \frac{s}{1 - s}\), \; \(\tilde{d} = \varphi \circ d\)
    \begin{itemize}
        \item \(\forall_{r > 0} \;\, B_d(x, r) = \{y \in X : \varphi(d(x, y)) < \varphi(r)\} = B_{\tilde{d}}\big(x, \frac{r}{1 + r}\big)\)
        \item \(\forall_{\varepsilon \in (0, 1)} \;\, B_{\tilde{d}}(x, \varepsilon) = B_d\big(x, \frac{\varepsilon}{1 - \varepsilon}\big)\) \, (wystarczy \(\varepsilon < 1\), bo dla \(\varepsilon \geq 1\) \; \(B_{\tilde{d}}(x, \varepsilon) = X\), a mniejsza kula też leży w \(U\))
    \end{itemize}
    zatem w warunku \(\forall_{x \in U} \; \exists_{r > 0} \;\, B(x, r) \subseteq U\) każdą kulę \(B_d\) można zastąpić kulą \(B_{\tilde{d}}\) o tym samym środku i odwrotnie \(\implies \mathcal{T}_d = \mathcal{T}_{\tilde{d}}\)
\end{podpunkty}

\bigskip
\zad{6} \; \((X, \mathcal{T})\) - przestrzeń topologiczna, \; \(A, B \subseteq X\)
\begin{podpunkty}
    \item \(\operatorname{int}(A \cap B) \overset{?}{=} \operatorname{int} A \cap \operatorname{int} B\)
    \begin{itemize}
        \item \(\subseteq\): \, \(x \in \operatorname{int}(A \cap B) \iff x \in \bigcup \{U \in \mathcal{T} : U \subseteq A \cap B\} \iff \exists_{U \in \mathcal{T}} \;\, x \in U \subseteq A \cap B \iff \exists_{U \in \mathcal{T}} \;\, x \in U \subseteq A \land x \in U \subseteq B \implies x \in \operatorname{int} A \cap \operatorname{int} B\)
        \item \(\supseteq\): \, \(x \in \operatorname{int} A \cap \operatorname{int} B \iff \exists_{U, V \in \mathcal{T}} \;\, x \in U \subseteq A \land x \in V \subseteq B \implies x \in U \cap V \overset{\text{iii)}}{\in} \mathcal{T} \land U \cap V \subseteq A \cap B \implies x \in \operatorname{int}(A \cap B)\)
    \end{itemize}
    \item \(\operatorname{int}(A \cup B) \overset{?}{=} \operatorname{int} A \cup \operatorname{int} B\)

    \smallskip
    kontrprzykład: \, \((\mathbb{R}, \mathcal{T}_{\text{st}})\), \; \(A = \mathbb{Q}, \; B = \mathbb{R} \setminus \mathbb{Q}\), \; \(\operatorname{int} \mathbb{R} = \mathbb{R} \neq \underbrace{\operatorname{int} \mathbb{Q}}_{= \, \emptyset} \cup \underbrace{\operatorname{int}(\mathbb{R} \setminus \mathbb{Q})}_{= \, \emptyset}\)
\end{podpunkty}

\bigskip
\zad{7} \; \(X\) - zbiór, \; \(\operatorname{CL}: 2^X \to 2^X\) \, t.że: \, (operacja domknięcia opisana aksjomatycznie - aksjomaty Kuratowskiego)
\begin{enumerate}[label=\roman*), leftmargin=1.9em, itemsep=0.2em, topsep=0.2em, parsep=0pt, partopsep=0pt]
    \item \(\operatorname{CL}(\emptyset) = \emptyset\)
    \item \(\operatorname{CL}(A) \supseteq A\)
    \item \(\operatorname{CL}(A \cup B) = \operatorname{CL}(A) \cup \operatorname{CL}(B)\)
    \item \(\operatorname{CL}(\operatorname{CL}(A)) = \operatorname{CL}(A)\)
\end{enumerate}
czy \(\mathcal{T} = \{U \subseteq X : \operatorname{CL}(U^c) = U^c\}\) - topologia na \(X\)?

\bigskip
\bigskip
\indeks{ć4}


\bigskip
\bigskip
\indeks{ć5}


\end{document}
